\chapter{Ideally exact categories}

\fontsize{12}{14.5}\selectfont

\section{Definition and characterization}

We are now ready to prove the theorem, which will lead to the definition of the ideally exact categories. The following characterization and subsequent results are based on the work of Janelidze \cite{Janelidze2024}.

\begin{mainbox}{}
    \begin{theo}\label{Ideally exact categories}
        Let $\mathcal{A}$ be a category. The following are equivalent:
        \begin{enumerate}
            \item $\mathcal{A}$ is exact and protomodular, has finite coproducts and the morphism $!:0 \to 1$ is a regular epimorphism.
            \item $\mathcal{A}$ is exact, has finite coproducts and there exists $\mathcal{X}$ semiabelian category with a monadic functor $U: \mathcal{A} \to \mathcal{X}$.
            \item There exist $\mathcal{X}$ semiabelian category with a monadic functor $U: \mathcal{A} \to \mathcal{X}$ and $F$ left adjoint to $U$ such that the associated monad $T=UF$ preserves regular epimorphisms and kernel pairs.
            \[
% https://tikzcd.yichuanshen.de/#N4Igdg9gJgpgziAXAbVABwnAlgFyxMJZABgBpiBdUkANwEMAbAVxiRAB12BbOnACwDGjYAEEAviDGl0mXPkIoyAJiq1GLNpx78hDYAA0JUmdjwEiARnKr6zVog7deg4YYB6wAKoAxI6phQAObwRKAAZgBOEFxIZCA4EEhWanZsniDUcHxYYThIALQWxiCR0UnUCbGZ2bkFybYaDt6S0iVRMYhxlYhK1A32jgIEgRkgfDB0UEhgTAwMFXRYDGyQYKzFpR298YmIyVk5eYiFG+3lO0jbB7XHRRRiQA
\begin{tikzcd}
    \mathcal{A} \arrow[dd, "U", shift left] \arrow[r, "\cong"]      & \mathcal{X}^{UF} \arrow[ldd, shift left] \\
                                                                    &                                          \\
    \mathcal{X} \arrow[uu, "F", shift left] \arrow[ruu, shift left] &                                         
    \end{tikzcd}
            \]
        \end{enumerate}
    \end{theo}
\end{mainbox}

\begin{mainbox}{}
    \begin{defi}
        A category $\mathcal{A}$ is said to be \textbf{ideally exact} if it satisfies one of the equivalent conditions of \Cref{Ideally exact categories}.
    \end{defi}
\end{mainbox}

\begin{subbox}{}
    \begin{proof}[\textit{Proof of \Cref{Ideally exact categories}}]\label{The proof}
        $\mathit{(1\implies 3)}$ Suppose $! : 0 \to 1$ is a regular epi. Since $\mathcal{A}$ is exact, the pullback functor $(!)^*: (\mathcal{A} \downarrow 1) \to (\mathcal{A} \downarrow 0)$ is monadic by the main result of the previous chapter: \Cref{monadic p^*}. Note that $(\mathcal{A} \downarrow 1) \cong \mathcal{A}$. Now consider $\mathcal{X} = (\mathcal{A} \downarrow 0)$, we want to prove that $\mathcal{X}$ is semiabelian. 
        
        $(\mathcal{A} \downarrow 0)$ is pointed, in fact if $(A,\alpha)$ is an object in $(\mathcal{A} \downarrow 0)$, we have the following initial and terminal objects, which coincide:
        \[
        % https://tikzcd.yichuanshen.de/#N4Igdg9gJgpgziAXAbVABwnAlgFyxMJZABgBpiBdUkANwEMAbAVxiRGJAF9T1Nd9CKAIzkqtRizYBBLjxAZseAkRFCx9Zq0TtZvRQKIBmUdQ2TtM7nv7KUxtaYladV+XyWDkAFhPjNbDk4xGCgAc3giUAAzACcIAFskERAcCCQAJkd-bQAdHMY0AAs6XRBYhKQyFLTEZLNnAEIAfUs5csTEY2qkAFYs8xA8guLS9srqVIzXMcQ+7sQfPwGhJsC2uI6uyYX+5yGGIpKgziA
        \begin{tikzcd}
        0 \arrow[r, "!_A"] \arrow[rd] & A \arrow[d, "\alpha"] &  & A \arrow[r, "\alpha"] \arrow[d, "\alpha"] & 0 \arrow[ld, "1_0"] \\
                                      & 0                     &  & 0                                         &                    
        \end{tikzcd}
        \]

        $\mathcal{A}$ has finite coproducts, so $(\mathcal{A} \downarrow 0)$ has finite coproducts. In fact consider $(A,\alpha)$, $(B,\beta)$ in $(\mathcal{A} \downarrow 0)$ and $A+B$ in $\mathcal{A}$, with the canonical injections $\iota_1 : A \to A+B$ and $\iota_2 : B \to A+B$. Let $(C,\gamma)$ be an object in $(\mathcal{A}\downarrow 0)$ with $f:A \to C$ and $g:B \to C$ morphisms in $(\mathcal{A}\downarrow 0)$ . Providing to $A+B$ the morphism $[\alpha,\beta]$ we have that $(A+B, [\alpha,\beta])$ is the coproduct of $(A,\alpha)$ and $(B,\beta)$ in $(\mathcal{A} \downarrow 0)$, in fact $\gamma [f,g] = [\alpha,\beta]$, so $[f,g]$ is also a morphism in $(\mathcal{A}\downarrow 0)$.
        \[
% https://tikzcd.yichuanshen.de/#N4Igdg9gJgpgziAXAbVABwnAlgFyxMJZAJgBoAGAXVJADcBDAGwFcYkQBBAagCEQBfUuky58hFGWLU6TVuwDCAoSAzY8BImQCM0hizaIQ5JcLVii5CrtkHOJlSPXjkAFis09cw337SYUAHN4IlAAMwAnCABbJEsQHAgkYkEwyJjELRoEpI8bdgAdfID6KKj6EBpGegAjGEYABUdzQ3CsAIALHHsI6KQAZizExDjPW0L8HHoAfS1utKQ3eKGRvMNxiEmp4gqQKtqGpo0Wts653sQBpaRMkFqwKCQAWj6V-XZQnb26xrMjkFaOl0UiAeulFtkMjQ7g8Lq8vCAAmd0pcIWQZG81vkmGh2uVKjVvodxP8TkDlKCFoMcuj4YVapMkbEqYg0aN2MhCtjcaQ6TBJpRGcNmTdoU8XABOXIYkDIUKkAIC-H7H6iP4A06+fhAA
\begin{tikzcd}
    A \arrow[rr, "\iota_1"] \arrow[rrdd, "f"', bend right] \arrow[rrd, "\alpha"'] &  & A+B \arrow[d] \arrow[d, "{[\alpha,\beta]}"] &  & B \arrow[ll, "\iota_2"'] \arrow[lldd, "g", bend left] \arrow[lld, "\beta"] \\
                                                                                  &  & 0                                                                                 &  &                                                                            \\
                                                                                  &  & C \arrow[u, "\gamma"']  \arrow[from=uu, "{[f,g]}"', bend right=49, crossing over, pos=0.54]                                                          &  &                                                                           
    \end{tikzcd}
    \]

    Now we want to prove that $(\mathcal{A} \downarrow 0)$ is exact, from the exactness of $\mathcal{A}$.
    First we prove it is finitely complete. It suffices to show that it has terminal a object and pullbacks. The terminal object is the one just shown above. It also has pullbacks, by \Cref{U creates con lim}, since pullbacks are connected limits and $\mathcal{A}$ has them.

    Consider now an equivalence relation $(R,r_1,r_2)$ in $(\mathcal{A} \downarrow 0)$, endowed with the morphism $\rho$ and $\alpha$ as shown in the diagram below. We want to prove that such equivalence relation is the kernel pair of some arrow. Since $\mathcal{A}$ is exact, it is a kernel pair in $\mathcal{A}$. By \Cref{equiv rel is kernel pair of the coeq} we can suppose that it is the kernel pair of its coequalizer.
    \[
% https://tikzcd.yichuanshen.de/#N4Igdg9gJgpgziAXAbVABwnAlgFyxMJZABgBpiBdUkANwEMAbAVxiRACUQBfU9TXfIRQBGclVqMWbAILdeIDNjwEiAJjHV6zVohAAhOXyWCio4eK1TdxbuJhQA5vCKgAZgCcIAWyRkQOCCRREDgACyxXHCDNSR0QdwB9VRBqBjoAIxgGAAV+ZSF4rAdQqJ43Tx9EPwDokPDIpABaYMs4xOFDEA9vWprEdQltNgBHFJBQmDooNhwAdwgJqYQyrorfaj6AZhih3QAdPfdQwNSMrNzjFV13IpLO7srgrZ2rEAPGNFC6e7X+jcDENsQAwsGA4lA6GF7GNWmwDjAAB5YOA4OAAQgABAdMjhvlwKFwgA
\begin{tikzcd}
    R \arrow[r, "r_2"', shift right] \arrow[r, "r_1", shift left] \arrow[rd, "\rho"'] & A \arrow[r, "q", two heads] \arrow[d, "\alpha"] & B \arrow[ld, "\exists! \beta", dashed] \\
                                                                                      & 0                                               &                                       
    \end{tikzcd}
    \]
    Since $r_1$ and $r_2$ are morphism in $(\mathcal{A} \downarrow 0)$, we have that $\alpha r_1 = \rho = \alpha r_2$. Thus, by the universal property of the coequalizer, there exists a unique morphism $\beta: B \to 0$ such that $\beta q = \alpha$. Providing to $B$ the morphism $\beta$, we have that $(B,\beta)$ is the coequalizer of $(R,r_1,r_2)$ in $(\mathcal{A} \downarrow 0)$.

    We proceed by proving that $(\mathcal{A} \downarrow 0)$ is regular. Both the conditions of the definition of regular category are easily extendable from $\mathcal{A}$ to $(\mathcal{A} \downarrow 0)$. In fact since in $\mathcal{A}$ each morphism admits a (regular epi, mono) factorization, then the same holds in $(\mathcal{A} \downarrow 0)$. Let $f=me$ be such a factorization through $Q$, with $f$ morphism in $(\mathcal{A} \downarrow 0)$. Endow $Q$ with the morphism $t$ induced by the universal property of the coequalizer. Now we want to prove that $e$ and $m$ are morphism in $(\mathcal{A} \downarrow 0)$. Then they will be respectively regular epi and mono in $(\mathcal{A} \downarrow 0)$. Notice first that $e$ is obviously a morphism in $(\mathcal{A} \downarrow 0)$, by the condition of $t$. Moreover, we have that $te=\alpha=\beta f= \beta me$ and since $e$ is a regular epi, hence by \Cref{reg epi is an epi} an epi, it follows that $m$ is a morphism in $(\mathcal{A} \downarrow 0)$.
    \[
    % https://tikzcd.yichuanshen.de/#N4Igdg9gJgpgziAXAbVABwnAlgFyxMJZABgBoBGAXVJADcBDAGwFcYkQBBEAX1PU1z5CKAEwVqdJq3YAhHnxAZseAkXKliEhizaIQARXn9lQtaRFapukMR4SYUAObwioAGYAnCAFskZEDgQSOqSOuxuRiCePn40gUhiIAAWMPRQ7DgA7hApaQg02tJ6bLzuXr6IifGIITj0WIwZ9Y0FVuy+pVHlsQFBiADMrWF6ADojTGhJ9JHRFbV9g6FFIGMARjB1M92VcQs0jFhg1lD0cCnpQ8tjMAAeWHA4cAAEAIRPOHbcQA
\begin{tikzcd}
    & Q \arrow[rd, "m", tail] \arrow[dd, "\exists ! t", pos=0.75, dashed] &                       \\
A \arrow[rr, "f", pos=0.25, crossing over] \arrow[ru, "e", two heads] \arrow[rd, "\alpha"] &                                                           & B \arrow[ld, "\beta"] \\
    & 0                                                         &                      
\end{tikzcd}
    \]

    In order to prove that regular epis are pullback stable in $(\mathcal{A} \downarrow 0)$ we can use the fact that they are in $\mathcal{A}$. In fact, we just endow the regular epi with the morphism induced by universal property of the coequalizer.

    We must now establish that the $(\mathcal{A} \downarrow 0)$ is protomodular. Since the construction of the change-of-base functor involves pullbacks, which are connected limits, then by \Cref{connected limits computed on the base} the construction is the same in $(\mathcal{A} \downarrow 0)$ and in $\mathcal{A}$. Thus $(\mathcal{A} \downarrow 0)$ inherits protomodularity from $\mathcal{A}$, noting that the definition of iso in $\mathcal{A}$ and $(\mathcal{A} \downarrow 0)$ is essentially the same, except the morphism condition. Hence, since $\mathcal{X} = (\mathcal{A} \downarrow 0)$ is exact, protomodular, pointed and has binary coproducts, we can conclude that it is semiabelian.

    It remains to prove that the monad $(!)^*(!)_!$ preserves regular epimorphisms and kernel pairs. Consider a regular epi $q:A \to B$ in $(\mathcal{A} \downarrow 0)$, it is the coequalizer of some parallel arrows. Since the coequalizer is a colimit, it is preserved by the left adjoint functor $(!)_!$. So it suffices to prove that, with $\mathcal{A}$ regular, $(!)^*$ preserves regular epis. Consider the following diagram:
    \[% https://tikzcd.yichuanshen.de/#N4Igdg9gJgpgziAXAbVABwnAlgFyxMJZABgBoBGAXVJADcBDAGwFcYkQBBEAX1PU1z5CKAEwVqdJq3YAhHnxAZseAkXKkREhizaIQ5ef2VCiZYlqm6QxADo28AW3gB9cl15HBq0aXM1t0nq29lhOcK5y3BIwUADm8ESgAGYAThAOSGQgOBBI6pI67ACOhiCp6Zk0OUgiHmVpGYj51Yi1CuWNAMxVuYjEdR1IACw9eQMNSN3ZvSMFgSAAFACEAJQAegBUC0UrPJTcQA
    \begin{tikzcd}
    0\times_1A \arrow[d] \arrow[rr, "(!)^*(q)"] &   & 0\times_1B \arrow[d] \\
    A \arrow[rr, "q"] \arrow[rd]                &   & B \arrow[ld]         \\
                                                & 1 &                     
    \end{tikzcd}
    \]
    By \Cref{p^*(f) is a pullback} the square is a pullback. Since $\mathcal{A}$ is regular, and hence so is $(\mathcal{A} \downarrow 0)$, it follows that $(!)^*(q)$ is a regular epi due to the stability of regular epis under pullbacks.

    Finally, we want to establish that the monad $(!)^*(!)_!$ preserves kernel pairs. This is obvious: $(!)_!$ preserves connected limits, such as kernel pairs, by \Cref{p_! preserves connected limits}, while $(!)^*$ preserves kernel pairs since it is a right adjoint. Thus $(!)^*(!)_!$ preserves kernel pairs.

    $\mathit{(3 \implies 2)}$ Suppose there exists $\mathcal{X}$ semiabelian category and a monadic functor $U: \mathcal{A} \to \mathcal{X}$. Assume that the associated monad $T$ preserves regular epimorphisms and kernel pairs.
    
    First we need to verify that $\mathcal{A}$ is exact. Since $U$ is monadic, we have that $\mathcal{A} \cong \mathcal{X}^T$. So it suffices to show that $\mathcal{X}^T$ is exact. We know that $\mathcal{X}$ is semiabelian, so it is exact. Then, by \Cref{algebras category is exact if the base category is exact}, $\mathcal{X}^T$ is exact.

    In order to prove that $\mathcal{A}$ has finite coproducts, it is first necessary to establish that $\mathcal{X}^T$ is protomodular. This follows directly from \Cref{algebras category is proto if the base category is proto} as $\mathcal{X}$ is protomodular. Since any protomodular category is Mal'tsev, then $\mathcal{X}^T$ is Mal'tsev. We already prove that $\mathcal{X}^T$ is also exact, then by \Cref{exact and Maltsev cat has reflexive coeq}, $\mathcal{X}^T$ has reflexive coequalizers. Since $\mathcal{X}$ has finite coproducts and $\mathcal{X}^T$ has reflexive coequalizers, then by \Cref{algebras category has colim if the base category has coprod}, $\mathcal{X}^T$ has finite coproducts. Finally, since $\mathcal{A} \cong \mathcal{X}^T$, it follows that $\mathcal{A}$  has finite coproducts.

    $\mathit{(2 \implies 1)}$ Suppose that $\mathcal{A}$ is exact, has finite coproducts and there exists a semiabelian category $\mathcal{X}$ and a monadic functor $U: \mathcal{A} \to \mathcal{X}$.

    Since $\mathcal{X}$ is protomodular, so is $\mathcal{X}^T$ by \Cref{algebras category is proto if the base category is proto}. Again, by monadicity, $\mathcal{A}$ inherits this properties.

    We are left to show that the morphism $!:0 \to 1$ is a regular epi. Since $\mathcal{A}$ is exact by assumption, the morphism $!$ admits a (regular epi, mono) factorization $! = me$.
    
    Now we aim to prove that if a morphism $m$ is a mono, then $U(m)$ acts as a mono in the morphisms in $\mathcal{X}$ between object coming from $\mathcal{A}$, ie if $f,g : U(A) \to U(B)$ for some $A$, $B$ in $\mathcal{A}$, then $mf=mg \implies f=g$ \footnote{This will be crucial to utilize the fully functor property of the comparison functor $K$.}. First, we prove that if $m$ is a mono, then $U^T(m)$ is a mono. Let $m$ be a mono and $A$ its domain. Then $(A, 1_A, 1_A)$ is the kernel pair of $m$, in fact, given $f$ and $g$ morphisms such that $mf=mg$, we obtain the universal property by the mono.
    \[
% https://tikzcd.yichuanshen.de/#N4Igdg9gJgpgziAXAbVABwnAlgFyxMJZARgBoAmAXVJADcBDAGwFcYkQBBEAX1PU1z5CKAMyli1Ok1bsuvftjwEiYqjQYs2iEAEUefEBkVCiZCeulbO+hYOUoADKQeSNM7QCEekmFADm8ESgAGYAThAAtkhiIDgQSE5SmuzEAPpyBmGR0TRxSGRJ7iBpXDSM9ABGMIwACgJKwiChWH4AFjg2IFlRiIl5iOQ0rTD0UEhgzIyMufRYjOw4s-MWydpR8l3hPQX9g7FLCwcrReuZW0gALLnxvTRVYGOIALQiiW5Wfp3dl9f5dzAPaJvSzsYJfc6IK6xG4xRhYMBWKD0ODDMbcSjcIA
\begin{tikzcd}
B \arrow[rdd, "g", bend right] \arrow[rrrd, "f", bend left] \arrow[rd, dashed] &                                       &  &                        \\
                                                                               & A \arrow[d, "1_A"] \arrow[rr, "1_A"'] &  & A \arrow[d, "m", tail] \\
                                                                               & A \arrow[rr, "m", tail]               &  & Q                     
\end{tikzcd}
    \]
    Since $U^T$ is a right adjoint, it preserves limit, such as kernel pairs. Thus also the following diagram is a kernel pair.
    \[
% https://tikzcd.yichuanshen.de/#N4Igdg9gJgpgziAXAbVABwnAlgFyxMJZABgBoBGAXVJADcBDAGwFcYkQBVAPQBUAKAIIBKEAF9S6TLnyEUAJlLFqdJq3bd+wsRJAZseAkQVUaDFm0Qhy2yfplEyS06oudegkaOUwoAc3hEoABmAE4QALZIAMw0OBBIZCrm7OQA+sAaHqI2IKER0bHxiOTOyZZpGe7C2TSM9ABGMIwAClIGsiCMMEE4OXmRiIlxSAogABYw9FBIYMyMjLH0WIzskGBspWqWHHzhnjr9SCUgw4ijOEsrlmsbSVucu56UokA
\begin{tikzcd}
U^T(A) \arrow[d, "1_{U^T(A)}"] \arrow[rr, "1_{U^T(A)}"] &  & U^T(A) \arrow[d, "U(m)"] \\
U^T(A) \arrow[rr, "U(m)"]                               &  & Q                       
\end{tikzcd}
    \]
    This implies that $U^T(m)$ is a mono by the universal property of the kernel pair. 
    
    It remains to prove that the image of a mono through the comparison functor $K : \mathcal{A} \to \mathcal{X}^T$ acts as a mono in the way described above. Then, since $U=U^TK$, we have that $U$ does the same. By monadicity, the functor $K$ is fully faithful, this is sufficient to prove that $K$ preserves monos. In fact, let $m$ be a mono in $\mathcal{A}$ and let $f$, $g$ be morphisms in $\mathcal{X}^T$ between object coming from $\mathcal{A}$ such that $K(m)f = K(m)g$. Let $f'$ and $g'$ be the morphisms in $\mathcal{A}$ such that $K(f') = f$ and $K(g') = g$ by the fully property of $K$. Then we have that
    \begin{align*}
    &K(m)f = K(m)g\\[1.5ex]
    &\implies K(m) K(f') = K(m) K(g')\\[1.5ex]
    &\implies K(mf') = K(mg')\\[1.5ex]
    &\implies mf' = mg'\\[1.5ex]
    &\implies f' = g'\\[1.5ex]
    &\implies K(f') = K(f')\\[1.5ex]
    &\implies f = g
    \end{align*}
    Thus $K(m)$ acts as a mono in the morphisms in $\mathcal{X}^T$ between object coming from $\mathcal{A}$. By the above discussion, $U(m)$ acts as a mono in the morphisms in $\mathcal{X}$ between object coming from $\mathcal{A}$.

    Picking up where we left off, consider $m$ the mono arising from the (regular epi, mono) factorization of $!:0 \to 1$ in $\mathcal{A}$. By the discussion above, $U(m)$ acts as a mono in the morphisms in $\mathcal{X}$ between object coming from $\mathcal{A}$. In particular, since the final object $1$ is a limit, it is preserved by $U$, then $U(m)$ acts as mono and has codomain $1$ in $\mathcal{X}$. Moreover, $\mathcal{X}$ is semiabelian, hence pointed, so $0=1$ in $\mathcal{X}$. This means that $U(m)$ acts as a mono and and has codomain $0$ in $\mathcal{X}$. Finally, this implies that $U(m)$ is an iso, in fact, suppose that $U(m)$ acts as a mono with codomain $0$ and let $U(C)$ be its domain. Denote by $!_{U(C)}$ the unique morphism from $U(C)$ to $0$, then $!_{U(C)}$ is the inverse of $U(m)$. In fact $!_{U(C)}U(m)=!_0=1_0$. Then composing with $U(m)$ we obtain $U(m)!_{U(C)}U(m)=U(m)$, thus, observing that $U(m)!_{U(C)}$ and $1_{U(C)}$ are a morphisms between objects arising from $\mathcal{A}$, we also have that $U(m)!_{U(C)}=1_{U(C)}$. Thus $U(m)$ is an iso.
    \[
% https://tikzcd.yichuanshen.de/#N4Igdg9gJgpgziAXAbVABwnAlgFyxMJZARgBoAGAXVJADcBDAGwFcYkRyQBfU9TXfIRTkK1Ok1bsAqgAoAwgEpuvEBmx4CRAEyiaDFm0QdlfdYKIBmXeIPT5SrmJhQA5vCKgAZgCcIAWyQyEBwIJBEbSSNZPwcVH38kHWDQxCsIwxAAQgB9YFlFLm5KLiA
\begin{tikzcd}
U(C) \arrow[r, "U(m)"] & 0 & 0 \arrow[r, "!_{U(C)}"] & U(C)
\end{tikzcd}
    \]
    
    Now we want to prove that $U$ reflects iso. As done before we prove that for the comparison functor $K$ and for $U^T$, then it follows for $U$. $K$ reflect iso because, by monadicity, it is fully faithful. In fact, if $K(f):K(X) \to K(Y)$ is an iso and $g: K(Y) \to K(X)$ is its inverse in $\mathcal{X}^T$, then $K(f)g=1_{K(Y)}$ and $gK(f)=1_{K(X)}$. Let $g'$ be the morphism in $\mathcal{A}$ such that $K(g') = g$ by the fully property of $K$. So we have that
    \begin{align*}
        &1_{K(Y)} = K(f)g\\[1.5ex]
        &\implies 1_{K(Y)} = K(f)K(g')\\[1.5ex]
        &\implies K(1_Y) = K(f)K(g')\\[1.5ex]
        &\implies K(1_Y) = K(fg')\\[1.5ex]
        &\implies 1_Y = fg'
    \end{align*}
    and similarly $1_X = g'f$. Thus $f$ is an iso in $\mathcal{A}$.

    Also $U^T$ reflects iso. In fact, let $U^T(k):U^T(V) \to U^T(W)$ be an iso in $\mathcal{X}^T$ and $h:U^T(W) \to U^T(V)$ its inverse. We aim to prove that $h$ is also a morphism in $\mathcal{X}^T$, so that it is an inverse for $k$ in $\mathcal{X}^T$. Let $\xi_V$ and $\xi_W$ be respectively the algebra structures of $U^T(V)$ and $U^T(W)$. Then, since $k$ is a morphism in $\mathcal{X}^T$, we have that $\xi_W T(k) = k \xi_V$, denoting $U^T(k)$ just as $k$. Composing with $h$ and $T(h)$ we obtain
    \[
    % https://tikzcd.yichuanshen.de/#N4Igdg9gJgpgziAXAbVABwnAlgFyxMJZABgBoBGAXVJADcBDAGwFcYkQA1EAX1PU1z5CKAEwVqdJq3YB1HnxAZseAkTLEJDFm0QgAKgAoOASnn9lQomI00t03YZmnuEmFADm8IqABmAJwgAWyQxEBwIJGJeXwDgxABmGnCkcmiQfyDIpIjEclspHRAAaxAaOAALLB8cJABaVIUMuNDkhPztdkMi0zLK6rqGmMy2sJzQiqqaxHr2+30DcudG2JTsrJAJ-um8yQ6HBedKbiA
\begin{tikzcd}
T(V) \arrow[d] \arrow[rr, "T(k)", shift left] &  & T(W) \arrow[d] \arrow[ll, "T(h)", shift left] \\
V \arrow[rr, "k", shift left]                 &  & W \arrow[ll, "T(h)", shift left]             
\end{tikzcd}
    \]
    \begin{align*}
        &\xi_W T(k) = k \xi_V\\[1.5ex]
        &\implies h \xi_W T(k) T(h)) = h k \xi_V T(h)\\[1.5ex]
        &\implies h \xi_W = \xi_V T(h)\\[1.5ex]
    \end{align*}
    Thus $h$ is a morphism in $\mathcal{X}^T$, and $k$ is an iso in $\mathcal{X}^T$. This implies, as discussed above, that also $U$ reflects iso.
    
    Finally we want to prove that $!:0 \to 1$ is a regular epi in $\mathcal{A}$. We already know that $U(m)$ is an iso. Since $U$ reflects iso, then $m$ is an iso. Thus $! = me$ is a regular epi.
    
    This holds because an iso, as $m$, is also a strong epimorphism, and $e$ is a regular epimorphism, hence it is also a strong epimorphism; therefore, the composition $me$ is a strong epimorphism. Since the $\mathcal{A}$ is regular, then by a standard result $!=me$ is a regular epi, as claimed.
    
\end{proof}
\end{subbox}

\begin{mainbox}{}
\begin{prop}
    An ideally exact category is finitely cocomplete.
\end{prop}
\end{mainbox}

\begin{subbox}{}
    \begin{proof}
        An ideally exact category is exact and protomodular, thus Mal'tsev. So by \Cref{exact and Maltsev cat has reflexive coeq} it has reflexive coequalizers. Then \Cref{algebras category has colim if the base category has coprod} implies that the category has all the finite colimits.
    \end{proof}
\end{subbox}

\begin{mainbox}{}
\begin{prop}\label{ideally exact has essentially nullary monad}
    Let $\mathcal{A}$ be an ideally exact category.
    There exist $\mathcal{X}$ semiabelian category with a monadic functor $U: \mathcal{A} \to \mathcal{X}$ and the corresponding monad is essentially nullary.
\end{prop}
\end{mainbox}

\begin{subbox}{}
    \begin{proof}
        We can consider $(!)^*$ to be such functor, as shown in the proof of \Cref{Ideally exact categories}$\mathit{:(1\implies 3)}$. Then by \Cref{p^* is essentially nullary} we have that the monad associated $(!)^*(!)_!$ is essentially nullary.
    \end{proof}
\end{subbox}

\section{The morphism-induced scenario}

As illustrated in \Cref{Ideally exact categories}$\mathit{:(1\implies 3)}$, in the definition of an ideally exact category, the monadic functor involved can be taken to be the pullback functor induced by the unique morphism $!:0 \to 1$. However, it may be more convenient to consider in the definition a different monadic functor $U$ satisfying certain properties, as in \Cref{Ideally exact categories}\textit{:(2),(3)}. That said, of course the monadic functor associated to $!:0 \to 1$ is always there. A natural question arises: is this monadic functor $U$, up to equivalence of categories, the same as the pullback functor associated to the unique morphism $!:0 \to 1$? In this section, we aim to address this question.

Given an adjunction $F \dashv U$ between categories finitely complete, with initial and terminal objects, and $\mathcal{X}$ pointed
\[
% https://tikzcd.yichuanshen.de/#N4Igdg9gJgpgziAXAbVABwnAlgFyxMJZABgBpiBdUkANwEMAbAVxiRAB12BbOnACwDGjYAEEAviDGl0mXPkIoAjOSq1GLNpx78hDYAA0JY1TCgBzeEVAAzAE4QuSMiBwQkykACMYYKE+r0zKyIIACqINRwfFjWOEgAtIpSNvaOiB6u-l4+fojOUTFxiIkB6sEgAGKSFGJAA
\begin{tikzcd}
\mathcal{A} \arrow[r, "U", shift left] & \mathcal{X} \arrow[l, "F", shift left]
\end{tikzcd}
\]
we can consider the adjunction induced on the initial objects $F^0 \dashv U^0$, defined as follows. Let $X$ be an object in $\mathcal{X}$ and let $(A,\alpha)$ be an object in $(\mathcal{A}\downarrow 0)$. Since $\mathcal{X}$ is pointed, then $\mathcal{X} \cong (\mathcal{X}\downarrow 0)$. The category of the initial objects $0$ in $\mathcal{X}$ and $\mathcal{A}$ is implicit. Denote $!^X:X \to 0$ the unique morphism in the zero object.
\begin{align*}
F^0(X) &= (F(X), F(!^X)) \\[1.5ex]
U^0(A,\alpha) &= 0\times_{UF(0)}U(A) = 0\times_{U(0)}U(A) = \mathrm{Ker}(U(\alpha))
\end{align*}
\[
% https://tikzcd.yichuanshen.de/#N4Igdg9gJgpgziAXAbVABwnAlgFyxMJZABgBoBGAXVJADcBDAGwFcYkRiQBfU9TXfIRQAmUsWp0mrdgFUAFAEEAlN14gM2PASKiqNBizaIQ84ip59NgomXH6pRjgB0neALbwA+sFNKu85W4JGCgAc3giUAAzACcINyQAZhocCCRye0N2Fzd6HAALGLdgNDiAKy5PYVVouITEMhBUpFFJLOMAQm9fLhqQWPj0lLTEVoNpY3kXJjR8+nM1AfrkppHG8cccvMLi0ogKz3IgriA
\begin{tikzcd}
0\times_{U(0)}U(A) \arrow[rr, "\mathrm{proj}_2"] \arrow[d, "\mathrm{proj}_1"] &  & U(A) \arrow[d, "U(\alpha)"] \\
0 \arrow[rr, "!_{U(0)}"]                                                      &  & U(0)                       
\end{tikzcd}
\]
Since $F$ is a left adjoint, it preserves colimits. So we can assume $F(0_\mathcal{X})=0_\mathcal{A}$.
\begin{equation}\label{induced ajunction on 0}
% https://tikzcd.yichuanshen.de/#N4Igdg9gJgpgziAXAbVABwnAlgFyxMJZABgBpiBdUkANwEMAbAVxiRAAoAdTgWzpwAWAY0bAAggF8ABNygQA7mDoAnZQqnEAlFIC8Urr37DRkmZzmKVa+VIBi7LZpATS6TLnyEUARnJVajCxsBnyCIgzAABrSsgpKqupaumahxhHRzv4wUADm8ESgAGZqPEhkIDgQSL4gAEYwYFBl1PTMrIggAKog1HACWIU4SAC03i5FJdXUlc11DU2I5X0DQ4ijLYHtILaZEkA
\begin{tikzcd}
(\mathcal{A} \downarrow 0) \cong (\mathcal{A} \downarrow F(0)) \arrow[r, "U^0", shift left] & (\mathcal{X} \downarrow 0) \cong \mathcal{X} \arrow[l, "F^0", shift left]
\end{tikzcd}
\end{equation}
One can prove that this is an adjunction. Given $\eta$ and $\epsilon$ the unit and counit of $F \dashv U$, then the unit of $F^0 \dashv U^0$ in an object $X$ in $\mathcal{X}$ is $\eta^0_X$
\begin{equation}\label{unit of adjunction in 0}
% https://tikzcd.yichuanshen.de/#N4Igdg9gJgpgziAXAbVABwnAlgFyxMJZARgBpiBdUkANwEMAbAVxiRAAYAdTvAW3gD6wAKoAKdgEoAvmIBiogBoSJIKaXSZc+QigDM5KrUYs2cxctXqQGbHgJF9AJkP1mrRCGHzJljbe1EZM7UriYe7L7WmnY6yOyk7C7G7iAKqoYwUADm8ESgAGYAThC8SPogOBBIjiHJbNwwOHQC7AC8AIRCXuLSINQMdABGMAwACtEBHoVYWQAWOJFFJUjxFVWIZEZu9Zy8dDizhbzAaMUAVgKOUovFpRvUldW12x7d7QB6SjfLiKuPiOVQiluHsDkcTucpAJiN87gAWB7rTbDMBQMqrIE7RrNNJqAq3JAItZlagotGIAC0ugxdQ8HzS-SGI3G-nsUxm81hhMRK36WDAKSgdDgs0yfS2YRADSa73YAlxFCkQA
\begin{tikzcd}
X \arrow[rrrd, "\eta_X", bend left] \arrow[rdd, "!^X"', bend right] \arrow[rd, "\eta^0_X", dashed] &                                                                                  &  &                              \\
                                                                                                   & 0\times_{UF(0)}UF(X) \arrow[rr, "\mathrm{proj_2}"] \arrow[d, "\mathrm{proj}_1=!^{U^0F^0(X)}"] &  & U(F(X)) \arrow[d, "UF(!^X)"] \\
                                                                                                   & 0 \arrow[rr, "\eta_0=!_{UF(0)}"']                                                &  & UF(0)                       
\end{tikzcd}
\end{equation}
where the outer diagram commutes because of the naturality of $\eta$, and the counit of $F^0 \dashv U^0$ in an object $(A,\alpha)$ in $(\mathcal{A} \downarrow 0)$ is
\[
\epsilon^0_{(A,\alpha)} = \epsilon_{A} F(\mathrm{proj}_2)
\]
\[
% https://tikzcd.yichuanshen.de/#N4Igdg9gJgpgziAXAbVABwnAlgFyxMJZABgBpiBdUkANwEMAbAVxiRADEAKAVQD1jOAQVIAdEYzQALOgEoZAXi7ExeALbwA+sG5KZAX25C5IPaXSZc+QigCM5KrUYs27Q4JkmzIDNjwEiAEz21PTMrIgggp7mvlZEdjYOoc4RuiYOMFAA5vBEoABmAE4Qqkh2IDgQSEGOYWxiMGjYDAQaUaYFxaWIZBVViOXJ4RycYqp0OJKFqsBoxQBWehoBHh0gRSVIvZVIAMwhTsNcYxNTM3MQi8sKXACEvMBcfALCYhLS+h7UDHQARjAMAAKFj81hAhSwWUkOGi6y61WoO0Q+1qKRAbwYUjo6T0QA
\begin{tikzcd}
{F(U^0(A,\alpha))=F(0\times_{UF(0)}U(A))} \arrow[r, "F(\mathrm{proj}_2)"] \arrow[rd, "{F(\mathrm{proj}_1)=F(!^{F(U^0(A,\alpha)})}"'] & FU(A) \arrow[r, "\epsilon_A"] & A \arrow[ld, "\alpha"] \\
                                                                                                                                     & F(0)                          &                       
\end{tikzcd}
\]
One can prove that the counit is a morphism in $(\mathcal{A} \downarrow 0)$ and that both $\eta^0$ and $\epsilon^0$ are natural transformations.

\begin{mainbox}{}
\begin{prop}\label{ajoint equivalence for pullback functor}
Given an adjunction $F \dashv U$
\[
% https://tikzcd.yichuanshen.de/#N4Igdg9gJgpgziAXAbVABwnAlgFyxMJZABgBpiBdUkANwEMAbAVxiRAB12BbOnACwDGjYAEEAviDGl0mXPkIoAjOSq1GLNpx78hDYAA0JY1TCgBzeEVAAzAE4QuSMiBwQkykACMYYKE+r0zKyIIACqINRwfFjWOEgAtIpSNvaOiB6u-l4+fojOUTFxiIkB6sEgAGKSFGJAA
\begin{tikzcd}
\mathcal{A} \arrow[r, "U", shift left] & \mathcal{X} \arrow[l, "F", shift left]
\end{tikzcd}
\]
in which $\mathcal{X}$ is pointed and protomodular, $U$ reflects isomorphisms, and the unit of the adjunction $\eta$ is cartesian.
Then the adjunction induced in $0$ $F^0 \dashv U^0$ is an adjoint equivalence.
\end{prop}
\end{mainbox}

\begin{subbox}{}
    \begin{proof}
        Consider the induced adjunction  $F^0 \dashv U^0$ in \eqref{induced ajunction on 0}. Observe that $(\mathcal{A} \downarrow 0) \cong \mathrm{Pt}_\mathcal{A}(0)$ via the unique morphism from the initial object.

        First, we aim to prove that $U^0$ reflects isos. Let $f:(A,\alpha)\to (A',\alpha')$ be a morphism in $\mathrm{Pt}_\mathcal{A}(0)$. Observe that $U^0(f)=(!_{U(0)})^*(U(f))$. 
        \[
% https://tikzcd.yichuanshen.de/#N4Igdg9gJgpgziAXAbVABwnAlgFyxMJZABgBpiBdUkANwEMAbAVxiRGIB0O8BbeAfWABVABTEAlAF9RAQXEhJpdJlz5CKAIzkqtRizaz5i5djwEiAZm3V6zVohCyA5EaUgMptUQAs13XbZObiw+OEFRCUkZJwU3D1VzFAAmUg0dW30HCNcTBPVkFKT0vXt2BR0YKABzeCJQADMAJwgeJDIQHAgkAFZqACMYMCgkAFoLdozSrh46HAALRp5gNGaAK0l+DRBqODmsepxRjWMQJpakLQ6uxF8QXf3DxBHLyYMRLkY0Oboc0+bWxApK5IW6vLLvDifb4ubYgBh0AYMAAKKjM6hAjSwVTmhx2ewOFxOZwB7U6FxsJTY01mCyWKwg634SVh8MRKM8iQxWJxsQa-3JwMBFIC4Pqv2JSCsgqBYJA1Pmi2Waw2SRiRP5iFJ1ylAyGkvaDCwYFKUDou0qsNlQgAesQRGKALwiACE4TEUnE1oAVCJRGLxRreoLQZSHK7hO7JCyETBkaivA4GDACerzogpWTEEHdcN0xNQ3KODMFXTlZsYtRWbH2Xk2Jjsbi7vjHsc3BKbtRM5d7gSni8C6JXXJeX807dM0Cey3hZlHC7BNFJPJKzG4xz0fWeamAUHM+0c-q8Q8jtuep3tf1Brmxu0pwKq2vaw5N4dJBRJEA
\begin{tikzcd}
0\times_{U(0)}U(A) \arrow[rrdd, "\mathrm{proj}_1", bend right, shift left] \arrow[r, "\mathrm{proj}_2"'] \arrow[rrrr, "U^0(f)=(!_{U(0)})^*(U(f))", dashed, bend left] & U(A) \arrow[rd, "U(\alpha)", shift left] \arrow[rr, "U(f)"] &                                                                                                   & U(A') \arrow[ld, "U(\alpha')"', shift right] & 0\times_{U(0)}A' \arrow[l, "\mathrm{proj}_2'"] \arrow[lldd, "\mathrm{proj}_1'"', bend left, shift right] \\
                                                                                                                                                                      &                                                             & U(0) \arrow[lu, "U(!_A)", shift left] \arrow[ru, "U(!_{A'})"', shift right]                       &                                              &                                                                                                          \\
                                                                                                                                                                      &                                                             & 0 \arrow[u, "!_{U(0)}"] \arrow[lluu, bend left, shift left] \arrow[rruu, bend right, shift right] &                                              &                                                                                                         
\end{tikzcd}
        \]
        Since $\mathcal{A}$ is protomodular, then $U(f)$ is an iso. By assumption $U$ reflects iso, thus $f$ is an iso, as claimed. 

        In order to prove that $F^0 \dashv U^0$ is an adjoint equivalence we need to prove that its unit $\eta^0$ and its counit $\epsilon^0$ are pointwise isos. By assumption $\eta$ is cartesian, thus, referring to \eqref{unit of adjunction in 0}, for each $X$ in $\mathcal{X}$ $\eta^0_X$ is an iso. Then by the triangular identity, for each $(A,\alpha)$ in $(\mathcal{A} \downarrow 0)$, we have that
        \[
        U^0(\epsilon^0_{(A,\alpha)}) \eta^0_{U^0((A,\alpha))} =\mathrm{Id}_{U^0(A,\alpha)}
        \]
        Hence $U^0(\epsilon^0_{(A,\alpha)})$ is an iso. We already proved that $U^0$ reflects iso, thus for each $(A,\alpha)$ in $(\mathcal{A} \downarrow 0)$ also $\epsilon^0_{(A,\alpha)}$ is an iso.
    \end{proof}
\end{subbox}

\begin{mainbox}{}
\begin{theo}
Let $\mathcal{A}$ be an ideally exact category, and let $U:\mathcal{A} \to \mathcal{X}$ be the monadic functor associated to the definition, coming with a left adjoint $F$. The functor $U$ is, up to category equivalence, the pullback functor $!^*$ if and only if the unit associated to $U$ is cartesian.
\begin{equation}\label{U is !^* up to iso}
% https://tikzcd.yichuanshen.de/#N4Igdg9gJgpgziAXAbVABwnAlgFyxMJZAJgBoAGAXVJADcBDAGwFcYkQAdDgW3pwAsAxk2ABBAL4hxpdJlz5CKMsWp0mrdl14DhjYAA1J02djwEi5UipoMWbRCAAUWvkJESABFygQA7mHoAJ0C-D3IASikZEAxTBQsKVVsNB2ceV10xcS8OH38gkN8PAEZI8VUYKABzeCJQADMQ7iRLEBwIJGKbdXsQAFUQGjh+LHqcJABaYuMQRohmxABmGnaWmn4Yeih2SDA2brtNDkECKqiGpqRlto7EMjVDhwBCAD0AKkGQYdHxxCmZuYLe6rRBdEAbLY7Aj7B4pTjHU7nWaXUErW6tb5jSZg5K9ABiSMBSGBt2umN+UwOcKeAH0nlJKOIgA
\begin{tikzcd}
(\mathcal{A} \downarrow 1) \arrow[rr, "\cong", no head] \arrow[dd, "!^*", shift left] &  & \mathcal{A} \arrow[dd, "U", shift left] \\
                                                                                      &  &                                         \\
(\mathcal{A} \downarrow 0) \arrow[rr, "\cong", no head] \arrow[uu, "!_!", shift left] &  & \mathcal{X} \arrow[uu, "F", shift left]
\end{tikzcd}
\end{equation}
\end{theo}
\end{mainbox}

\begin{subbox}{}
    \begin{proof}
        We already proved that $(\mathcal{A} \downarrow 1) \cong \mathcal{A}$, by the unique morphism to the terminal object. Suppose that the unit $\eta$ associated to $U$ is cartesian, we aim prove that $(\mathcal{A} \downarrow 0) \cong \mathcal{X}$ and that the diagram commutes up to natural isomorphisms.

        By the definition of ideally exact category we have that $\mathcal{X}$ is semiabelian, thus $\mathcal{X}$ is pointed and protomodular. Moreover, since $U$ is monadic, $U$ reflects isos. By assumption we also have that $\eta$ is cartesian. Then, by \Cref{ajoint equivalence for pullback functor} the adjunction $F^0 \dashv U^0$ in \eqref{induced ajunction on 0} is an adjoint equivalence, thus an equivalence of categories.
        
        We need to prove that such equivalence of categories makes the diagram commute up to natural iso.
        \[
        % https://tikzcd.yichuanshen.de/#N4Igdg9gJgpgziAXAbVABwnAlgFyxMJZAJgBoAGAXVJADcBDAGwFcYkQAdDgW3pwAsAxk2ABBAL4hxpdJlz5CKMsWp0mrdl14DhjYAA1J02djwEi5UipoMWbRCAAUWvkJESABFygQA7mHoAJ0C-D3IASikZEAxTBQsKVVsNB2ceV10xcS8OH38gkN8PAEZI8VUYKABzeCJQADMQ7iRLEBwIJGKbdXsQAFUQGjh+LHqcJABaYuMQRohmxABmGnaWmn4Yeih2SDA2brtNDkECKqiGpqRlto7EMjVDhwBCAD0AKkGQYdHxxCmZuYLLo3NZfEZjSbA5K9ABi51mlzuK1u12+EL+UJ67CeAH0nvDAUh7qtEMC0b9MY9+i9yJ9GPQAEYwRgABTkZkUIECWCq-HGAMRwJJ93JnQOKRAMJpdMZzLZcXMDm5vP5lHEQA
\begin{tikzcd}
(\mathcal{A} \downarrow 1) \arrow[rr, "\cong", no head] \arrow[dd, "!^*", shift left]    &  & \mathcal{A} \arrow[dd, "U", shift left]                                 \\
                                                                                         &  &                                                                         \\
(\mathcal{A} \downarrow 0) \arrow[uu, "!_!", shift left] \arrow[rr, "U^0"', shift right] &  & \mathcal{X} \arrow[uu, "F", shift left] \arrow[ll, "F^0"', shift right]
\end{tikzcd}
        \]
        The upward-left directed square commutes up to natural iso by definition of $F^0$. Then, up to category equivalence, $!_!\cong F$. Since $F \dashv U$ and $!_! \dashv \ !^*$, then, since the right adjoint is unique up to natural iso, $!^*\cong U$. Hence, also the downward-right directed square commutes up to natural iso.
        
        Conversely, suppose that $U$ is, up to natural iso, the pullback functor of $!$, then, by \Cref{eta cartesian}, the unit is cartesian.
    \end{proof}
\end{subbox}

\begin{rembox}{}
    \begin{rem}
        One can also prove directly that downward-right directed square \eqref{U is !^* up to iso} commutes up to natural iso. So, we aim to prove that $U^0 !^*$ corresponds to the functor $U$. In order to do that we compute $U^0 !^*$ in an object $A$ in $\mathcal{A}$. Consider the object $(A,!^A)$, and its pullback
        \[
        % https://tikzcd.yichuanshen.de/#N4Igdg9gJgpgziAXAbVABwnAlgFyxMJZABgBpiBdUkANwEMAbAVxiRGIB0O8BbeAfQCMAQRABfUuky58hFACZyVWoxZtREqdjwEiZQcvrNWiduMkgM22UUUHqRtacHjlMKAHN4RUADMAThA8SILUOBBIAMwOqiYgAIQAehoWAUFIiiDhUTHGbPHmfoHBiGRZEYihKnmmXDx0OAAW-jzAaIEAVmL88oUgaSVl2YiZjnF1Dc2t7RBdQq5iQA
\begin{tikzcd}
0\times_1A \arrow[rr, "\mathrm{proj}_2"] \arrow[d, "\mathrm{proj}_1"] &  & A \arrow[d, "!^A"] \\
0 \arrow[rr, "!"]                                                     &  & 1                 
\end{tikzcd}
        \]
        Since $!^*$ is a right adjoint, also the following square is a pullback.
        \[
        % https://tikzcd.yichuanshen.de/#N4Igdg9gJgpgziAXAbVABwnAlgFyxMJZABgBpiBdUkANwEMAbAVxiRAFUAKYgSgB0+eALbwA+gEYuAQR4gAvqXSZc+QigBM5KrUYs2U+YpAZseAkTLjt9Zq0QdushUtOqimq9Rt774+dpgoAHN4IlAAMwAnCCEkcWocCCQAZi9dOwcAQgA9GUMI6NjETRBElLTbNi5MpyMomKQyUqTEeJ1K+y4BITocAAtIoWA0aIArOVF1WoKGxCay4oqfB27egaGRiHGJJwo5IA
\begin{tikzcd}
U(0)\times_1U(A) \arrow[rr, "U(\mathrm{proj}_2)"] \arrow[d, "U(\mathrm{proj}_1)"] &  & A \arrow[d, "U(!^A)"] \\
U(0) \arrow[rr, "U(!)"]                                                           &  & 1                    
\end{tikzcd}
\]
Then, apply $U^0$ to $(U(0)\times_1U(A),U(\mathrm{proj}_1))$. By definition \eqref{induced ajunction on 0}, we need to compute the pullback of $U(\mathrm{proj}_1)$ along $!_U(0)$.
Since the right square is a pullback, as we just proved, and the whole diagram is a pullback, then the left square is the pullback of $U(\mathrm{proj}_1)$ along $!_U(0)$. Thus the result of $U^0 !^*$ in an object $A$ in $\mathcal{A}$ is the pullback $0\times_1U(A)$.
\[
% https://tikzcd.yichuanshen.de/#N4Igdg9gJgpgziAXAbVABwnAlgFyxMJZAJgBoAGAXVJADcBDAGwFcYkQBVACnIEoAdfngC28APoBGbgEFeIAL6l0mXPkIoALBWp0mrdtIVKQGbHgJEyEnQxZtEnHnMXKzaolus1b+hxKOuqhYo5KReunbs5AEmKubqyKFU3nr2IOSCIuJSXLIKOjBQAObwRKAAZgBOEMJIEjQ4EEgAzCmRDtwAhAB6eS4gVTVIZCCNLW2+jp3OxoO1iKGjTYj1EZPcgsL0OAAWlcLAaNUAVvJixDMV1fOLY4gjPmkb-Fu7+4cnZxKXA9dIWkthhM0p0xMBuHx5DE5kgAKwNZbkfowxDwwGIDTySjyIA
\begin{tikzcd}
0\times_1U(A) \arrow[rr] \arrow[d] &  & U(0)\times_1U(A) \arrow[rr, "U(\mathrm{proj}_2)"] \arrow[d, "U(\mathrm{proj}_1)"] &  & A \arrow[d, "U(!^A)"] \\
0 \arrow[rr, "!_{U(0)}"]           &  & U(0) \arrow[rr, "U(!)"]                                                           &  & 1                    
\end{tikzcd}
\]
Finally we prove that the pullback $0\times_1U(A)$ is just $U(A)$. Observe that the following square is the pullback of $U(!^A)$ along $!$. In fact, given two morphism $f_1:X \to U(A)$, $f_2:X \to 0$ such that the external diagram commutes, there exists a unique morphism such the two triangles commutes. Such unique morphism is $f_2$.
\[
% https://tikzcd.yichuanshen.de/#N4Igdg9gJgpgziAXAbVABwnAlgFyxMJZABgBpiBdUkANwEMAbAVxiRAA0QBfU9TXfIRQBGUsKq1GLNgFUAFAEEAlN14gM2PASIBmMRPrNWiEPOWq+mwUVEAmA1OMhiAXmEX1-LUOR771Q2kTdy4JGCgAc3giUAAzACcIAFskPRAcCCQAFh44xJTEW2oM7NyQBOSkMnTMxFEQBiwwJyg6OAALcJAAxzZYgH0QtQqC6pLC6gAjGDAoVOrApwGhvMrEMdq06dmkAFodBd6TAdsPEaR68bTFtgBCAD1gMyUuM-yL4tqiySM2YX6nooXtwKFwgA
\begin{tikzcd}
X \arrow[rd, "f_1", dashed] \arrow[rrrd, "f_1", bend left] \arrow[rdd, "f_2", bend right] &                                                   &  &                \\
                                                                                          & U(A) \arrow[d, "!^{U(A)}"] \arrow[rr, "1_{U(A)}"] &  & U(A) \arrow[d] \\
                                                                                          & 0=1 \arrow[rr]                                    &  & 1             
\end{tikzcd}
\]
Thus the pullback $0\times_1U(A)$ of $U(!^A)$ along $!$ is just $A$. This means that also the downward-right directed square commutes up to natural isomorphisms.
\end{rem}
\end{rembox}

This previous theorem answer the question we posed at the beginning of this section. The following example shows that we already came across to such a scenario.

\begin{rembox}{}
    \begin{ex}
In the first chapter we saw the same theorem in the particular case of $\mathcal{A}=\mathsf{Ring}$ and $\mathcal{X}=\mathsf{Ring_*}$, where $1=\{0\}$ and $0=\mathbb{Z}$. 
\[
% https://tikzcd.yichuanshen.de/#N4Igdg9gJgpgziAXAbVABwnAlgFyxMJZABgBpiBdUkANwEMAbAVxiRAAoAdTgWzpwAWAIwBmwAEpYwAcwC+AAm5QIAdzB0AThtXyAjAEoQs0uky58hFGQBMVWoxZt5XXv2FjJMhUtXqtO7j5BISFgAC1ZQ2NTbDwCImtyO3pmVkQQQLdRCSk5IxMQDFiLBNJbahTHdMzgj1yAfQAqWSM7GChpeCJQEW0eJDIQHAgkXWo4ASwRHCQAWjH7VLYAQgA9RvyevtHqYYHxyem5hcq0kGX65c2QXoh+xEShkcQAZgOpmcR5ioczgFVrrd7m8nkhHhMPscfkt0gAxQHbRCDPYPaFVDKcADGBGkIF2dCwDDYkDArGoAhgdCgxIIrGiN0RCxRIIhR0QJ1+bBqAg0PGAAGsYBoWtQGHQhDAGAAFMxxSwgDRYaQCGb0oFIEEohasz4cmEcWGkbhoLD1bkhcKRPEgMUS6WykrpRXK1UUWRAA
\begin{tikzcd}
(\mathsf{Ring} \downarrow 1) \arrow[dd, "!^*", shift left] \arrow[rr, "\cong", no head]                       &  & \mathsf{Ring} \arrow[dd, "U", shift left]                                                    \\
                                                                                                              &  &                                                                                              \\
 (\mathsf{Ring} \downarrow \mathbb{Z}) \arrow[uu, "!_!", shift left] \arrow[rr, "\mathrm{Ker}"', shift right] &  & \mathsf{Ring_*} \arrow[uu, "F", shift left] \arrow[ll, "{(F,\pi_\mathbb{Z})}"', shift right]
\end{tikzcd}
\]
In particular the bottom category equivalence constructed in that case exactly corresponds to \eqref{induced ajunction on 0}
    \end{ex}
\end{rembox}

\section{Ideals and quotients}

In this section we will investigate the notion of ideals in an ideally exact category and their corresponding quotients. We aim to emulate the classical definition of ideals in a ring, and the corresponding quotients. As already mentioned, recall that a proper ideal in a ring is a non-unital ring. Note that the notion of normal subobjects is well defined in a semiabelian category as the kernel of a morphism. Also recall that a non-unital ring is semiabelian, while a unital ring fails to be, as mentioned in the first chapter. Given the previous observations, one can model the notion of ideals as the normal subobjects through a functor, which is the forgetful functor in rings, having a semiabelian category as codomain. Note that this framing is consistent with the characterization of ideally exact categories given in \Cref{Ideally exact categories}. In particular in $\mathsf{Ring}$ we have the following structure.
\[
% https://tikzcd.yichuanshen.de/#N4Igdg9gJgpgziAXAbVABwnAlgFyxMJZABgBpiBdUkANwEMAbAVxiRAB12BbOnACwBGAM2AAlLGADmAXxDTS6TLnyEUARnJVajFm049+wsRMkB9AFSzpWmFEnwioIQCcIXJGRA4ISDduasiCAAqnIU0kA
\begin{tikzcd}
\mathsf{Ring} \arrow[r, "U"] & \mathsf{Ring_*}
\end{tikzcd}
\]
In general in an ideally exact category $\mathcal{A}$ we have the following structure, with semiabelian $\mathcal{X}$.
\[
\begin{tikzcd}
\mathcal{A} \arrow[r, "U"] & \mathcal{X}
\end{tikzcd}
\]

We start with a technical lemma.

\begin{mainbox}{}
    \begin{lemma}\label{induced algebra between morphisms}
        Let $\mathcal{X}$ be a category with finite coproducts and let $T$ be an essentially nullary monad in $\mathcal{X}$. Given $(X_1,\xi_1)$ and $(X_2,\xi_2)$ $T$-algebras in $\mathcal{X}$, $X$ an object in $\mathcal{X}$, $m: X_1 \to X$ a morphism and $n: X \to X_2$ a mono, such that their composition is a $T$-algebra morphism, then there exists a unique $T$-algebra structure $\xi$ on $X$ making $m$ and $n$ $T$-algebra morphisms.
        \begin{equation}\label{diag induced algebra between morphisms}
% https://tikzcd.yichuanshen.de/#N4Igdg9gJgpgziAXAbVABwnAlgFyxMJZABgBoBGAXVJADcBDAGwFcYkQANAfXJAF9S6TLnyEU5CtTpNW7Dv0EgM2PASIAmSTQYs2iTl3UKhK0UTLEpO2foAqACm7kAlMaXDVY5JsvaZekAdudVcBExE1cVJfaV12INCpGCgAc3giUAAzACcIAFskMhAcCCQJWJsQArCQHPyymhKkTQqAwkb6LEZ2HE7umrqCxABmRtLEIusAgB1pgA8sHjdBpAAWMea-OP1ZhcNl3KHR4vGAVi3KhzzQxRXEc5O1i4CHMBusw6QHpsRyxiwwAEoPQ4AALZIgZ7sXZYfiUPhAA
\begin{tikzcd}
T(X_1) \arrow[d, "\xi_1"] \arrow[r, "T(m)"] & T(X) \arrow[r, "T(n)"] \arrow[d, "\xi", dashed] & T(X_2) \arrow[d, "\xi_2"] \\
X_1 \arrow[r, "m"]                          & X \arrow[r, "n", tail]                          & X_2                      
\end{tikzcd}
        \end{equation}
    \end{lemma}
\end{mainbox}
\begin{subbox}{}
    \begin{proof}
First we show that if the there exists a $T$-algebra structure $\xi$ making $n$ a $T$-algebra morphism, then necessarily $m$ is a $T$-algebra morphism. In fact, since $n$ is a mono, and both the outer and the right diagram in \eqref{diag induced algebra between morphisms} commute, then the left diagram commutes.
\begin{align*}
    \xi_2 & T(n) T(m) = n m \xi_1 \\[1.5ex]
    & \implies n \xi T(m)= n m \xi_1 \\[1.5ex] 
    & \implies \xi T(m) = m \xi_1
\end{align*}


Now we aim to prove that there exists a unique $T$-algebra structure $\xi$ on $X$ making $n$ a $T$-algebra morphism. Since $T$ is an essentially nullary monad, then $[T(!_{X}),\eta_{X}]$ is a strong epi. Consider the following diagram of solid arrows. 
\[
% https://tikzcd.yichuanshen.de/#N4Igdg9gJgpgziAXAbVABwnAlgFyxMJZAJgBoAGAXVJADcBDAGwFcYkQAVACgA0BKEAF9S6TLnyEU5CtTpNW7buT4BqHkJEgM2PASJlishizaIQPAPrENonRKLTDNYwrPrht8XpRkAjEflTTl4rAUFZGCgAc3giUAAzACcIAFskXxocCCRpORN2ZG4AQgtgHkE+UgAdKpgcelLyyhsQJNT0zOzEAGZnQIKUmoAPLAtfYsaxitJfCx5mmkZ6ACMYRgAFMV1JEESsKIALHBa2tMRcrKQAFj78s24wMM1T686kMjzXEGHR6w9W5JnXogS6ID71LCMdgQqG3L6ERYrNabOzeXb7I4nQE5N49RZYMBBKD0OAHSIgOFBH4UkBLVYbLb2Mx7Q7HcKCIA
\begin{tikzcd}
T(0)+X \arrow[rr, "{[T(!_{X}),\eta_{X}]}"] \arrow[dd, "{[m\xi_1T(!_{X_1}),1_X]}"'] &  & T(X) \arrow[d, "T(n)"] \arrow[lldd, "\xi"', dashed] \\
                                                                                   &  & T(X_2) \arrow[d, "\xi_2"]                           \\
X \arrow[rr, "n"', tail]                                                           &  & X_2                                                
\end{tikzcd}
\]
The following compositions with the canonical injections of $T(0)+X$ prove that the diagram commutes.
\begin{itemize}[itemsep=1.5ex]
    \item $\xi_2 T(n) T(!_X) = \xi_2 T(n!_X) = \xi_2 T(!_{X_2})$
    \item $nm\xi_1 T(!_{X_1}) = \xi_2 T(nm) T(!_{X_1}) = \xi_2 T(nm !_{X_1}) = \xi_2 T(!_{X_2})$
    \item $\xi_2T(n) \eta_X = \xi_2 \eta_{X_2} n = n$
    \item $n 1_X = n$
\end{itemize}
Since $[T(!_{X}),\eta_{X}]$ is a strong epi, there exists a unique morphism $\xi$ making the diagram commute. $\xi$ turns out to be a algebra for $X$. 

It remains to prove that it is unique. Suppose that there exists another algebra structure $\xi'$ on $X$ making $n$ a $T$-algebra morphism. Then
\begin{align*}
    n & \xi = \xi_2 T(n) = n \xi' \\[1.5ex]
    & \implies n \xi = n \xi' \\[1.5ex]
    & \implies \xi = \xi'
\end{align*}
    \end{proof}
\end{subbox}

\begin{mainbox}{}
\begin{cor}\label{T-algebra in the reflexive relation}
    Let $\mathcal{X}$ be a category with finite coproducts and let $T$ be an essentially nullary monad in $\mathcal{X}$. Given a $T$-algebra $(X,\xi)$ in $\mathcal{X}$, then every reflexive relation $R$ on $\mathcal{X}$
    \[
    % https://tikzcd.yichuanshen.de/#N4Igdg9gJgpgziAXAbVABwnAlgFyxMJZABgBpiBdUkANwEMAbAVxiRACUQBfU9TXfIRQBGclVqMWbABrdxMKAHN4RUADMAThAC2SMiBwQkokHAAWWNTiQBaAEzV6zVohAaA+sO68QmnXupDY2pzS2tEBwlnNg87EGoGOgAjGAYABX48AhisRTNrHnUtXUQTIMR9JylXBATk1IzsLKEQLDBsWDkuIA
\begin{tikzcd}
R \arrow[r, "r_1", shift left=2] \arrow[r, "r_2"', shift right=2] & X \arrow[l]
\end{tikzcd}
    \]
    admits a unique $T$-algebra structure making $R$ a reflexive relation on $\mathcal{X}^T$.
\end{cor}
\end{mainbox}

\begin{subbox}{}
    \begin{proof}
        Let $s$ be the common section of $r_1$ and $r_2$. Then
        \[\langle r_1, r_2 \rangle s = \langle 1_X, 1_X \rangle \]
        \[
        % https://tikzcd.yichuanshen.de/#N4Igdg9gJgpgziAXAbVABwnAlgFyxMJZARgBoAGAXVJADcBDAGwFcYkQAlEAX1PU1z5CKAEwVqdJq3YANADpy8AW3gACGTz4gM2PASLlxNBizaIQG7hJhQA5vCKgAZgCcISpGJA4ISQ5NN2BF5nNw9Efx8kMgDpcwVGejBbRhhVFwB9MnSMkVUFFySUthoceixGdjKKzVD3T1LfRBiAIxgwKCQAWgBmfxM4kASi1NViDJlScfk5QuTUkBpEtsYABQE9YRAXLFsACxweSm4gA
\begin{tikzcd}
X \arrow[r, "s"] \arrow[rr, "{\langle 1_X,1_X\rangle}"', bend right] & R \arrow[r, "{\langle r_1, r_2 \rangle}", tail] & X\times X
\end{tikzcd}
\]
Observe that ${\langle 1_X,1_X\rangle}$ is a $T$-algebra morphism, in fact given the product structure from \Cref{product lift} of $X\times X$ on the category $\mathcal{X}^T$
\[
% https://tikzcd.yichuanshen.de/#N4Igdg9gJgpgziAXAbVABwnAlgFyxMJZABgBpiBdUkANwEMAbAVxiRABUAKADQEoQAvqXSZc+QigBM5KrUYs2XbgB1leALbwABH0HCQGbHgJEALDOr1mrRBx78hIo+LOlJsqwtvc9TsSZQyd0t5GxAfRwNRYwlkaWC5azYVNSxNOB1fKOcAuNIARg9QxXtVDW0lB1kYKABzeCJQADMAJwh1JDIQHAgkc0SvEFUADyws1vakaW7exABmEKTbEbHIiY7EfOoepAA2RcHVBjowWoYYLS5VNCwAfXzeUkvOa7vJXlUWk7PWamOAIxgDAACtEXLZzk0cOM2ht9jMkABWA5hFZaMppbQrGGTTbbWbTTxhK7KG63d44jbIhHzFFsV7kylILY0rpEkoMh4gP50QEgsEBEAtLC1AAW0LWsKR+L6dOWpLu+W5IABQNBOQkQpF4sEFAEQA
\begin{tikzcd}
T(X) \arrow[dd, "\xi"] &  & T(X\times X) \arrow[d, "{\langle T(\pi_1), T(\pi_2)\rangle}"] \arrow[rr, "T(\pi_2)"] \arrow[ll, "T(\pi_1)"'] &  & T(X) \arrow[dd, "\xi"] \\
                       &  & T(X)\times T(X) \arrow[d, "\xi \times \xi"]                                                                  &  &                        \\
X                      &  & X\times X \arrow[rr, "\pi_2"] \arrow[ll, "\pi_1"']                                                           &  & X                     
\end{tikzcd}
\]
we have that the following diagram commutes, easily composing with the canonical projections $\pi_1$ and $\pi_2$.
\[
% https://tikzcd.yichuanshen.de/#N4Igdg9gJgpgziAXAbVABwnAlgFyxMJZABgBpiBdUkANwEMAbAVxiRABUAKADQEoQAvqXSZc+QigBM5KrUYs2XbgB1leALbwABH0HCQGbHgJFpARln1mrRBx69VG7Uv5CRR8UTKTL8myG49dzETKVIfaisFWxU1LE04HUFZGCgAc3giUAAzACcIdSQyEBwIJABmSL82VQAPLBBqBjoAIxgGAAVRYwkQBhhsnCCQPIKi6lKkMyrrRU5VZrA0-rMAfW5SNe4tVVy6Jf7+Jtb2ro9QvoGhtxH8wsRpkrLEaTlZ2wX95ZgtLlU0LCrMy8Ui-ebKAGrSQOZR7A6sY5tTrdTy2fqDYaje6vSaIAAsM2iIDqWB2cQSZPqjT6J2R5166Ou+ixFQmzwJbyJn3hW026zJcO+1OaSLOIQZV2SAiAA
\begin{tikzcd}
T(X) \arrow[dd, "\xi"] \arrow[rr, "{T(\langle1_X,1_X \rangle)}"] &  & T(X\times X) \arrow[d, "{\langle T(\pi_1), T(\pi_2)\rangle}"] \\
                                                                 &  & T(X)\times T(X) \arrow[d, "\xi \times \xi"]                   \\
X \arrow[rr, "{\langle1_X,1_X \rangle}"]                         &  & X\times X                                                    
\end{tikzcd}
\]
Thus we have $s$ and $\langle r_1,r_2 \rangle$ morphisms whose composition is the $T$-algebra morphism $\langle 1_X,1_X \rangle$. Then, by \Cref{induced algebra between morphisms}, there exists a unique $T$-algebra structure on $R$ making $s$ and $\langle r_1,r_2 \rangle$ $T$-algebra morphisms.
    \end{proof}
\end{subbox}

In the next results we consider the following lattice. Let $A$ be an object in a category $\mathcal{A}$.
\begin{itemize}[itemsep=1.5ex]
    \item $\mathrm{Quot}(A)$: regular epis, up to iso, with domain $A$;
    \item $\mathrm{ERel}(A)$: equivalence relations, up to iso, on $A$;
    \item $\mathrm{RRel}(A)$: reflexive relations, up to iso, on $A$.
\end{itemize}

Finally, if the category is semiabelian, we consider also the lattice consisting of normal subobjects $\mathrm{NSub}(A)$, containing the kernels, up to iso, of some morphisms with domain $A$.

The following theorem shows the correspondence between the lattices of quotients and normal subobjects in a semiabelian category. This is a generalization of the classical result in $\mathsf{Grp}$.

\begin{mainbox}{}
\begin{theo}\label{quotients in semiabelian}
    Let $\mathcal{X}$ be a semiabelian category. Then for each object $X$ in $\mathcal{X}$ the following lattices are isomorphic.
    \[\mathrm{Quot}(X) \cong \mathrm{NSub}(X)\]
\end{theo}
\end{mainbox}

\begin{subbox}{}
    \begin{proof}
    Consider the map 
    \[\mathrm{Quot}(X) \to \mathrm{NSub}(X)\]
    that associates to each regular epi $q:X \to Q$ its kernel
    \[
    % https://tikzcd.yichuanshen.de/#N4Igdg9gJgpgziAXAbVABwnAlgFyxMJZARgBoAGAXVJADcBDAGwFcYkQANEAX1PU1z5CKAEwVqdJq3YBFHnxAZseAkXLiaDFm0QgAOnoC29HAAsAToeABpGOe4AKAI4BKHhJhQA5vCKgAZuYQhkjqIDgQSGSS2uxOIDSmMPRQ7DgA7hBJKQi8AUEhiGLhkYjk3JTcQA
\begin{tikzcd}
\mathrm{Ker}(q) \arrow[r] & X \arrow[r, "q", two heads] & Q
\end{tikzcd}
\]
    and the map 
    \[\mathrm{NSub}(X) \to \mathrm{Quot}(X)\]
    that associates to each kernel pair of some morphism $f:X \to Y$ the regular epi $p$ underlying the corresponding (regular epi, mono) factorization of $f$.
    \[
% https://tikzcd.yichuanshen.de/#N4Igdg9gJgpgziAXAbVABwnAlgFyxMJZABgBoBGAXVJADcBDAGwFcYkQAdDgW3pwAsATt2ABpGIIC+ACgBmAShCTS6TLnyEU5CtTpNW7ABpKVIDNjwEiAZh00GLNohABNE6osaiAJlLFdDgbOAApKujBQAObwRKCyghDcSNogOBBIvnqO7LLuIPGJSGSp6YjkynEJSWU0aUi2IPww9FDsOADuEE0tCPb6TmZ5BdUNdYiZOPRYjG1TM5KUkkA
\begin{tikzcd}
                          &                                              & P \arrow[rd, tail] &   \\
\mathrm{Ker}(f) \arrow[r] & X \arrow[rr, "f"] \arrow[ru, "p", two heads] &                    & Y
\end{tikzcd}
    \]

    Consider a regular epi $q$, and its kernel. Then consider the factorization of $q$. Since $q$ is a regular epi, we can assume, up to iso, that the regular epi underlying the factorization is $q$ itself with the identity morphism as mono. Thus the composition
    \[\mathrm{Quot}(X) \to \mathrm{NSub}(X) \to \mathrm{Quot}(X)\] is the identity morphism. 
    
    Consider a kernel pair of some morphism $f:X \to Y$. Then consider the regular epi underlying the factorization of $f$. By \Cref{f and e share the kernel} the regular epi of the factorization and $f$ share the same kernel. Thus the composition
    \[\mathrm{NSub}(X) \to \mathrm{Quot}(X) \to \mathrm{NSub}(X)\]
    is the identity morphism.
    We proved that the two maps are inverses of each other, hence the isomorphism holds.
    \end{proof}
\end{subbox}

Now we aim to generalize the previous theorem to an ideally exact category.

\begin{mainbox}{}
\begin{theo}\label{quotients in ideally exact}
    Let $\mathcal{A}$ be an ideally exact category and let $U:\mathcal{A} \to \mathcal{X}$ be the underlying functor. Then for each object $A$ in $\mathcal{A}$ the following lattices are isomorphic.
    \[\mathrm{Quot}(A) \cong \mathrm{NSub}(U(A))\]
\end{theo}
\end{mainbox}

\begin{subbox}{}
    \begin{proof}
        Let $A$ be an object in $\mathcal{A}$. First we have that
        \[\mathrm{Quot}(A) \cong \mathrm{ERel}(A)\]
        in fact $\mathcal{A}$ is exact, hence kernel pairs and equivalence relations correspond to each others, and by \Cref{coeq of the kernel pair} and \Cref{kernel pair of the coeq} we have the given isomorphism. Furthermore
        \[\mathrm{ERel}(A) \cong \mathrm{RRel}(A)\]
        because $\mathcal{A}$ is protomodular, hence Mal'tsev. By \Cref{ideally exact has essentially nullary monad} the unit of the adjunction underlying $U$ is essentially nullary, then, by \Cref{T-algebra in the reflexive relation}, we have that
        \[\mathrm{RRel}(A)\cong \mathrm{RRel}(U^T(A)) \]
        Since $U$ is monadic we also have that
        \[\mathrm{RRel}(A)\cong \mathrm{RRel}(U(A)) \]
        By definition, the codomain of $U$ is semiabelian, thus exact and Mal'tsev, then as above 
        \[\mathrm{RRel}(U(A)) \cong \mathrm{ERel}(U(A)) \cong \mathrm{Quot}(U(A))\]
        Finally, by the previous result \Cref{quotients in semiabelian} we have that
        \[\mathrm{Quot}(U(A)) \cong \mathrm{NSub}(U(A))\]
    \end{proof}
\end{subbox}

\begin{rembox}{}
    \begin{rem}
        Note that in an ideally exact category $\mathcal{A}$, given an object $A$ in $\mathcal{A}$, the lattice $\mathrm{NSub}(A)$ is not well defined, since $\mathcal{A}$ is in general non-pointed.
    \end{rem}
\end{rembox}

\begin{rembox}{}
    \begin{rem}
        In $\mathsf{Ring}$, given an object $A$ in $\mathcal{A}$, the lattice $\mathrm{NSub}(U(A))$ consists of ideals. We have the following corresponding morphisms, composed by an ideal $I$, ie a kernel, and its corresponding quotient map, ie a regular epi, $\pi$.
        \[
        % https://tikzcd.yichuanshen.de/#N4Igdg9gJgpgziAXAbVABwnAlgFyxMJZARgBoAGAXVJADcBDAGwFcYkQAlEAX1PU1z5CKAEwVqdJq3YcA9AEkefEBmx4CRcuJoMWbRCEXcJMKAHN4RUADMAThAC2SLSBwQkZEHAAWWazmcdKX0QAB1QtCwlG3snRDFXd0RyY24gA
\begin{tikzcd}
I \arrow[r, hook] & R \arrow[r, "\pi"] & R/I
\end{tikzcd}
        \]
    \end{rem}
\end{rembox}

With this final theorem we proved that in an ideally exact category quotients and normal subobjects correspond to each others. This is exactly what happens in $\mathsf{Ring}$ as mentioned in the first chapter. In fact in $\mathsf{Ring}$ the ideals, which are the normal subobjects in $\mathsf{Ring_*}$ through the forgetful functor, correspond to quotients, which are the regular epis. Observe that in $\mathsf{Grp}$ the result holds directly from \Cref{quotients in semiabelian}, since $\mathsf{Grp}$ is a semiabelian category, and normal subgroups are objects of $\mathsf{Grp}$. Thus in this case we do not need to consider an underlying functor, or, equivalently, we can just consider the identity functor.